% Part: first-order-logic
% Chapter: model-theory
% Section: partial-iso

\documentclass[../../../include/open-logic-section]{subfiles}

\begin{document}

\olfileid[tr]{mod}{bas}{dlo}
\section{Yoğun Doğrusal Sıralamalar}

\begin{defn}
  \emph{Uç noktaları olmayan yoğun doğrusal sıralama}, tek bir iki yerli
  !!{predicate}~$<$ içeren !!{language} için aşağıdaki tümceleri
  sağlayan !!a{structure}~$\Struct M$ ifadesidir:
  \begin{enumerate}
  \item $\lforall[x][\lnot x < x]$;
  \item $\lforall[x][\lforall[y][\lforall[z][(x < y \lif (y < z \lif x
    <z ))]]]$;
  \item $\lforall[x][\lforall[y][(x< y \lor \eq[x][y] \lor y < x)]]$;
  \item $\lforall[x][\lexists[y][x < y]]$;
  \item $\lforall[x][\lexists[y][y < x]]$;
  \item $\lforall[x][\lforall[y][(x < y \lif \lexists[z][(x < z \land
        z < y))]]]$.
 \end{enumerate}
\end{defn}

\begin{thm}\ollabel{thm:cantorQ}
  Uç noktaları olmayan herhangi iki !!{enumerable} yoğun doğrusal
  sıralama izomorftur.
\end{thm}

\begin{proof}
  $\Struct M_1$ ve $\Struct M_2$, uç noktaları olmayan !!{enumerable}
  yoğun doğrusal sıralamalar; ${<_1}=\Assign{<}{M_1}$ ve
  ${<_2}=\Assign{<}{M_2}$ olsun. Aralarındaki bütün kısmi
  izomorfizmaların kümesine $\PIso I$ diyelim. En azından
  $\emptyset\in\PIso I$ olduğundan $\PIso I$ boş değildir. $\PIso I$
  kümesinin ileri-geri özelliğini sağladığını göstereceğiz. O zaman
  $\Struct M_1\iso[p]\Struct M_2$ olur ve teorem
  \olref[pis]{thm:p-isom1} sonucundan çıkar.

  $\PIso I$ kümesinin ileri özelliğini sağladığını göstermek için
  $p\in\PIso I$ ve $i=1,\dots,n$ için $p(a_i)=b_i$ olsun; genelliği
  bozmadan $a_1<_1a_2<_1\cdots<_1a_n$ varsayalım. Verilen
  $a\in\Domain{M_1}$ için $b\in\Domain{M_2}$ elemanını şöyle bulun:
  \begin{enumerate}
  \item $a<_1a_1$ ise $b<_2b_1$ olacak bir $b\in\Domain{M_2}$ alın;
  \item $a_n<_1a$ ise $b_n<_2b$ olacak bir $b\in\Domain{M_2}$ alın;
  \item bazı $i$ için $a_i<_1a<_1a_{i+1}$ ise
    $b_i<_2b<_2b_{i+1}$ olacak bir $b\in\Domain{M_2}$ alın.
  \end{enumerate}
  $\Struct M_2$ uç noktaları olmayan yoğun bir doğrusal sıralama
  olduğundan istenen özelliğe sahip bir $b$ her zaman bulunabilir.
  $q=p\cup\{\langle a,b\rangle\}$ tanımlayın; böylece $q\in\PIso I$,
  $p$'nin istenen genişlemesidir. Bu, ileri özelliğini gösterir. Geri
  özelliği benzerdir. Dolayısıyla $\Struct M_1\iso[p]\Struct M_2$;
  \olref[pis]{thm:p-isom1} gereği $\Struct M_1\iso\Struct M_2$.
\end{proof}

\begin{prob}
  $\PIso I$ kümesinin geri özelliğini sağladığını doğrulayarak
  \olref[mod][bas][dlo]{thm:cantorQ} teoreminin ispatını tamamlayın.
\end{prob}

\begin{rem}
  $\Struct S$, uç noktaları olmayan herhangi bir !!{enumerable} yoğun
  doğrusal sıralama olsun. O zaman \olref{thm:cantorQ} gereği
  $\Struct S\iso\Struct Q$; burada $\Struct Q=(\Rat,<)$, evreni
  rasyonel sayılar kümesi~$\Rat$ olan !!{enumerable} yoğun doğrusal
  sıralamadır. Şimdi \olref[thm]{remark:R} içindeki
  !!{structure}~$\Struct R=(\Real,<)$ ifadesini yeniden ele alalım.
  $\Struct R\elemequiv\Struct S$ olacak !!a{enumerable}
  !!{structure}~$\Struct S$ bulunduğunu görmüştük. Ancak $\Struct S$,
  uç noktaları olmayan !!a{enumerable} yoğun doğrusal sıralamadır;
  dolayısıyla !!{structure}~$\Struct Q$ ile izomorf (ve bu yüzden
  elementer eşdeğer) olur. Elementer eşdeğerliğin geçişliliği gereği
  $\Struct R\elemequiv\Struct Q$. (Aynı ileri-geri savıyla doğrudan
  $\Struct R\iso[p]\Struct Q$ göstererek de bu sonuca ulaşabilirdik.)
\end{rem}
\end{document}
